% new_version_main.tex

\documentclass{article}

\PassOptionsToPackage{numbers,compress}{natbib}
\usepackage{neurips_2026}

\usepackage[utf8]{inputenc}
\usepackage[T1]{fontenc}
\usepackage{hyperref}
\hypersetup{hidelinks,hypertexnames=false}
\usepackage{url}
\usepackage{booktabs}
\usepackage{amsmath}
\usepackage{amsfonts}
\usepackage{amssymb}
\usepackage{amsthm}
\usepackage{nicefrac}
\usepackage{microtype}
\usepackage{xcolor}
\usepackage{graphicx}
\usepackage{algorithm}
\usepackage{algpseudocode}

\newcommand{\R}{\mathbb{R}}
\newcommand{\E}{\mathbb{E}}
\newcommand{\1}{\mathbf{1}}
\newcommand{\argmin}{\operatorname*{arg\,min}}
\newcommand{\argmax}{\operatorname*{arg\,max}}
\newcommand{\ri}{\operatorname{ri}}
\newcommand{\diag}{\operatorname{Diag}}

\theoremstyle{plain}
\newtheorem{theorem}{Theorem}
\newtheorem{lemma}{Lemma}
\newtheorem{proposition}{Proposition}
\newtheorem{corollary}{Corollary}

\theoremstyle{definition}
\newtheorem{definition}{Definition}
\newtheorem{assumption}{Assumption}

\theoremstyle{remark}
\newtheorem{remark}{Remark}

\title{Decision-Focused Scenario Learning for Contextual Stochastic Programming}

\author{Anonymous Authors}

\begin{document}
\maketitle

\begin{abstract}
Decision-focused scenario learning for contextual two-stage stochastic programming has emerged as a promising research direction, but open questions remain. A general theory for when scenario generators can yield optimal decision policies is currently lacking. Many existing methods require strict assumptions on which second-stage data are allowed to be random. Generally, the non-differentiability of the decision-focused objective is a source of difficulty. This paper makes three theoretical contributions, each aiming to remedy one of these challenges. First, we present a comprehensive theory for when scenario generators can induce optimal policies, showing that learning a single out-of-support scenario component is often sufficient. Second, we demonstrate a theoretical and methodological equivalence result between learning different second-stage data components, relaxing assumptions required for applying existing methods. Third, we study log-barrier smoothing enabled training, presenting representability and optimization results. Finally, leveraging the theory, we propose a novel method inspired by interior-point methods from classical optimization, and validate it on standard benchmarks.
\end{abstract}

\section{Introduction}
\label{sec:introduction}

Machine learning is increasingly used to support high-stakes decisions in logistics, energy, healthcare, finance, and public planning \citep{rudin2019stop,bertsimas2020predictive,mandi2024decision}. Yet automated decision making is not new. Since the early development of mathematical programming, classically associated with Dantzig's work on linear programming \citep{dantzig1963linear}, computers have solved large constrained optimization problems according to deterministic rules. Today, mathematical programming is essential for the operation of power systems, supply chains, transportation networks, financial systems, and other critical infrastructure \citep{wood2013power,simchi_levi2014logic,ahuja1993network,markowitz1952portfolio}. Its ability to strictly handle complicated constraints has made it resistant to being replaced by pure machine learning approaches \citep{birge2011introduction}.

This has motivated growing interest in hybrid methods that combine machine learning with optimization. Rather than replacing mathematical programming, these methods use learning to improve the inputs, approximations, or solution procedures of optimization models \citep{bengio2021machine,kotary2021endtoend,mandi2024decision}. A central example is decision-focused learning (DFL), where a predictive model is trained through the downstream optimization problem. The model learns to output parameters that produce good decisions, not merely accurate predictions \citep{donti2017task,elmachtoub2022smart,wilder2019melding}.

We study this idea for contextual two-stage stochastic programming. In a two-stage stochastic program, a first-stage decision is made before uncertainty is realized, and a second-stage recourse decision is made afterward \citep{birge2011introduction,shapiro2014lectures}. For example, an energy operator may commit generation before demand and renewable output are known, and then dispatch reserves after uncertainty is observed. In the contextual setting, the decision maker observes side information, such as weather forecasts, market data, or recent demand, before choosing the first-stage decision. The aim of contextual stochastic programming is to learn a map from context data to high-quality decisions.

Following the recent survey of \citet{sadana2025survey}, existing approaches to contextual stochastic programming can be grouped into three broad categories.

\begin{itemize}
    \item \textbf{Predict-then-optimize.} Learn the conditional distribution, and then solve the resulting stochastic program. When the learned distribution is accurate, this can be combined with classical optimization to obtain an exact optimal policy for the estimated problem. However, distribution learning can be data intensive, and small statistical errors can cause large downstream decision errors \citep{elmachtoub2022smart,donti2017task}.

    \item \textbf{Direct policy learning.} Learn the map from context to decision. This can be effective when the optimal policy is simpler than the full conditional distribution. However, learned policies often struggle to satisfy complicated constraints \citep{kotary2021endtoend}.

    \item \textbf{Decision-focused learning.} Learn scenarios, costs, or other optimization inputs to optimize downstream decision quality. Because decisions are produced by solving an optimization problem, feasibility is naturally respected. The main challenge is that training requires repeatedly solving optimization problems and differentiating through their solution maps, which can quickly become computationally expensive \citep{amos2017optnet,agrawal2019differentiable,mandi2024decision}.
\end{itemize}

This paper focuses on the third category. Specifically, we study decision-focused learning of single synthetic scenarios. Given a context, a model outputs one scenario, a decision is computed, optimized for the scenario, and then evaluated on the true downstream objective. Single-scenario learning is attractive because both training and inference are much faster than solving large sample-average approximations, whose size grows with the number of scenarios \citep{shapiro2014lectures}. The challenge is that optimal-solution maps are often non-differentiable, making gradient-based training difficult \citep{wilder2019melding,mandi2020interior}.

We make three contributions.

\begin{itemize}
    \item \textbf{Optimal-policy realizability via single-scenario learning.}
    We prove that, under weak well-posedness assumptions, optimal stochastic-programming decisions can be obtained by solving carefully chosen single-scenario problems. In particular, full-scenario predictions are often redundant, and either pure cost or demand predictions are sufficient for representing optimal decisions.

    \item \textbf{Equivalence between learning cost and demand.}
    We show that under general conditions optimal decisions can be represented with demand forecasts or cost forecasts. For example, a problem with random demand and deterministic cost can be solved by producing a well-tuned cost forecast, just as well as a demand forecast. Which component one chooses can therefore be decided by computational considerations and feasibility, rather than being dictated by the structure of randomness in the original problem.

    \item \textbf{Log-barrier smoothing and trainability.}
    We study the theoretical properties of log-barrier smoothing as a means of enabling decision-focused scenario learning. We show that it induces smoothness in the scenario-parametrized loss, with closed-form derivatives. Additionally, we prove an unexpected generalized convexity property for the same loss function, showing that it is componentwise invex in both the cost and demand. The approach hence avoids a known theoretical pathology associated with quadratic smoothing. Motivated by these results, we propose a neural net training method inspired by classical interior point methods.

\end{itemize}

\paragraph{Closest approaches in the literature.}
Our work is closest to the random-cost decision-focused learning literature of \citet{elmachtoub2022smart} and \citet{wilder2019melding}, the log-barrier differentiable-optimization approach of \citet{mandi2020interior}, and decision-focused stochastic-optimization methods such as \citet{donti2017task} and \citet{islip2025contextual}. Prior stochastic-programming methods often learn distributions or multiple scenarios, and frequently rely on quadratic smoothing, differentiable relaxations, or black-box surrogates \citep{donti2017task,agrawal2019differentiable,islip2025contextual}. The closest single-scenario and decision-optimal forecasting results are due to \citet{estes2023smart} and \citet{homemdemello2024forecasting}, but these usually impose restrictions on which second-stage component is random. Our contribution is to give a realizability theory for scenario generation as a policy class, extend and connect existing methods, and prove novel theoretical properties for log-barrier smoothing.

\section{Notation and Problem Statement}
\label{sec:notation-problem}
\label{sec:problem}

\subsection{Data set and scenarios}
\label{sec:data-scenarios}

A context variable $x\in \R^{n_x}$ is represented as a Euclidean vector. A scenario is represented as a collection of three \emph{scenario components}
\[
\xi=\bigl(W(\xi),h(\xi),q(\xi)\bigr)
\in \R^{m_2\times n_2}\times\R^{m_2}\times\R^{n_2},
\]
where the component $W(\xi)\in \R^{m_2\times n_2}$ represents a second-stage constraint matrix (recourse matrix), the component $h(\xi)\in \R^{m_2}$ represents a second-stage constraint vector (demand vector), and the component $q(\xi)\in \R^{n_2}$ represents a second-stage cost vector (cost vector).

We assume the context and scenarios are sampled from a set $\mathcal X \times \Xi$ according to an unknown joint distribution $P$, and that we have access to a dataset $\mathcal D$ of $N$ context-scenario pairs
\begin{equation}
\label{eq:dataset}
\mathcal D:=\{(x_i,\xi_i)\}_{i=1}^N\subseteq \mathcal X\times \Xi,
\qquad
(x_i,\xi_i)\overset{\mathrm{i.i.d.}}{\sim}P.
\end{equation}

Throughout, $A\in\R^{m_1\times n_1}$, $b\in\R^{m_1}$, and $c\in\R^{n_1}$ denote first-stage data, and $T\in\R^{m_2\times n_1}$ denotes technology data, all assumed to be deterministic.

\subsection{Contextual stochastic programming with log-barrier smoothing}
\label{sec:contextual-logbarrier}

For each $x\in \mathcal X$, and each admissible choice of log-barrier smoothing parameter $\mu \in [0,\infty)$, the log-barrier smoothed two-stage stochastic linear program is given by
\begin{equation}
\label{eq:Pmu-x}
\begin{aligned}
(\mathcal{P}_\mu^{x})\qquad
&\min_{z\in Z}
\left\{
c^\top z-\mu\sum_{j=1}^{n_1}\log z_j
+
\E_{\boldsymbol{\xi}\sim P(\cdot\mid x)}\!\left[Q_\mu(z,\boldsymbol{\xi})\right]
\right\},
\\
&Z:=\{z\in\R^{n_1}_+:\;Az=b\}.
\end{aligned}
\end{equation}
where $Q_\mu$ denotes the second-stage pay-off
\begin{equation}
\label{eq:Qmu}
\begin{aligned}
Q_\mu(z,\xi)
&:=
\min_{y\in\mathcal Y(z,\xi)}
q(\xi)^\top y-\mu\sum_{k=1}^{n_2}\log y_k,
\\
\mathcal Y(z,\xi)
&:=
\left\{
y\in\R^{n_2}_+:\; W(\xi)y=h(\xi)-Tz
\right\}.
\end{aligned}
\end{equation}
We emphasize that our main goal is to solve the unsmoothed problem, defined by setting $\mu = 0$ in all of the above. We study the log-barrier smoothed problem primarily as a means to this end. For compactness, many of our theoretical results are stated for $\mu\geq 0$ to cover both cases simultaneously.

For our theoretical analysis, we use the following standing assumptions:
\begin{enumerate}
    \item \label{ass:finite-support} $\operatorname{supp}(\mathbb P)=\mathcal X\times\Xi$ is assumed to be compact. The conditional distribution $P(\cdot\mid x)$ is assumed to be well defined for all $x\in\mathcal X$. And its support $\Xi_x=\operatorname{supp}(\mathbb P(\cdot\mid x))$ is assumed to be finite.
    \item \label{ass:augmented-recourse} (Augmented relatively complete recourse) For every $\xi\in\Xi$ and every $z\in Z\cap\R^{n_1}_{++}$, the second-stage feasible set contains a strictly positive vector, $\mathcal Y(z,\xi)\cap\R^{n_2}_{++}\neq\emptyset$.
    \item \label{ass:no-redundant-second-stage-equalities} (No redundant second-stage equalities) The first-stage constraint matrix $A$ has full row rank. The same holds for the recourse matrix $W(\xi)$ for every scenario in the support $\xi\in\Xi$.
    \item \label{ass:logbarrier-well-posedness} (Log-barrier smoothed well-posedness) For every $x\in\mathcal X$, the problem $\mathcal{P}_\mu^{x}$ has a solution for every value of $\mu\geq 0$.
    \item \label{ass:fixed-technology-matrix} (Fixed technology matrix) General stochastic programs allow the second-stage technology matrix $T$ to depend on randomness. We restrict ourselves to cases where this does not hold.
\end{enumerate}
The first three assumptions are non-restrictive and standard in stochastic programming. The log-barrier well-posedness assumption may seem opaque, but any well-posed linear program can generally be reformulated to respect it. This is what allows interior-point methods to function as a go-to standard for solving large LPs. See, e.g., \citet{wright1997primal} for an overview of the classic theory.

Throughout, we will use the following equivalent formulation of $\mathcal{P}_\mu^{x}$, in the classical form of a stochastic optimization problem with a random objective:
\begin{equation}
\label{eq:Gmu-min}
(\mathcal{P}_\mu^{x})\qquad
\min_{z\in\R^{n_1}}
\E_{\boldsymbol{\xi}\sim P(\cdot\mid x)}\!\left[\hat G_\mu(z,\boldsymbol{\xi})\right].
\end{equation}
Where, $\hat G_\mu$ denotes the scenario-specific first-stage cost
\begin{equation}
\label{eq:hatGmu}
\hat G_\mu(z,\xi):=
c^\top z
-\mu\sum_{j=1}^{n_1}\log z_j
+Q_\mu(z,\xi)
+\iota_{Z}(z),
\end{equation}
where $\iota_{Z}(z)=0$ if $z\in Z$ and $\iota_{Z}(z)=+\infty$ otherwise.

For the original $\mu = 0$ problem, $\mathcal{P}_\mu^{x}$ may have multiple optimal solutions. However, for the smoothed problem with $\mu > 0$, the solution is unique. For strictly positive $\mu > 0$, we denote the unique optimal policy by
\begin{equation}
\label{eq:zsharp-mu}
z_\mu^\#:\mathcal X\to Z,
\qquad
z_\mu^\#(x):=\argmin_{z\in\R^{n_1}}G_\mu(z,x),
\end{equation}

We will frequently omit the $\mu$ subscript when considering the unsmoothed, $\mu = 0$ case. By classical theory, the problem returns to its original non-smoothed version as $\mu$ goes to zero:
\begin{equation}
\label{eq:G-limit}
G(z,x):=\lim_{\mu\downarrow 0}G_\mu(z,x),
\qquad
\hat G(z,\xi):=\lim_{\mu\downarrow 0}\hat G_\mu(z,\xi).
\end{equation}
In particular, we define
\begin{equation}
\label{eq:zsharp-limit}
z^\#(x)=\lim_{\mu \downarrow 0} z_\mu^\#(x) \in \argmin_{z\in Z} G(z,x)
\end{equation}
as the unique optimal policy over the unsmoothed problem, defined by the log-barrier smooth limit.

\subsection{Synthetic scenarios}
\label{sec:synthetic-scenarios}

Throughout, we will frequently consider ``synthetic'' feasible scenarios
\begin{equation}
\label{eq:Xi-synthetic}
\Xi_\mu^{\mathrm{synthetic}}:=
\left\{
(W,h,q)\in
\R^{m_2\times n_2}\times\R^{m_2}\times\R^{n_2}:
-\infty<
\min_{z\in Z}\hat G_\mu(z,(W,h,q))
<\infty
\right\}.
\end{equation}
Intuitively, $\Xi_\mu^{\mathrm{synthetic}}$ includes all parameters of second-stage data that yield a well-posed single-scenario problem. In particular, it includes all supported scenarios, $\Xi\subseteq \Xi_\mu^{\mathrm{synthetic}}$. Synthetic scenarios generally are not in the support $\Xi$, and may lack meaningful physical interpretation. In such cases they are best viewed as parameters for encoding decisions, not as interpretable forecasts.

The uniqueness of the log-barrier optimum allows us to define a decision map over synthetic scenarios for $\mu>0$:
\begin{equation}
\label{eq:single-scenario-map-mu}
\hat z_\mu^\#:\Xi_\mu^{\mathrm{synthetic}}\to Z,
\qquad
\hat z_\mu^\#(\xi):=\argmin_{z\in\R^{n_1}}\hat G_\mu(z,\xi),
\end{equation}
and the scenario-parametrized decision problem as
\begin{equation}
\label{eq:scenario-parametrized-decision-problem}
\min_{\xi\in\Xi_\mu^{\mathrm{synthetic}}}
\E_{\boldsymbol{\xi}'\sim P(\cdot\mid x)}
\!\left[
\hat G_\mu\!\left(\hat z_\mu^\#(\xi),\boldsymbol{\xi}'\right)
\right].
\end{equation}

\subsection{Decision focused scenario generation}
\label{sec:decision-focused-scenario-generation}

As a learning problem, our contextual stochastic optimization reduces to solving the empirical minimization problem defined by
\begin{equation}
\label{eq:direct-policy-erm}
\min_{\eta\in\Theta}\ \frac1N\sum_{i=1}^N \hat G_\mu(\pi_\eta(x_i),\xi_i),
\qquad
\pi_\eta:\mathcal X\to Z,
\end{equation}
where $\eta$ parametrizes a family of policies.
The decision focused scenario learning problem will be defined by:
\begin{equation}
\label{eq:scenario-generation-erm}
\min_{w\in\Omega}\ \frac1N\sum_{i=1}^N
\hat G_\mu\!\left(\hat z_\mu^\#(\phi_w(x_i)),\xi_i\right),
\qquad
\phi_w:\mathcal X\to \Xi_\mu^{\mathrm{synthetic}}.
\end{equation}
We remark that decision focused scenario generation can be viewed as a specific case of direct policy learning, by setting the policy family in \eqref{eq:direct-policy-erm} as equal to all policies obtained by composing $\hat z_\mu^{\#}$ with the policies in the scenario policy class:
\begin{equation}
\label{eq:scenario-policy-class}
\Pi_{\mathrm{SG}}:=
\left\{
\pi_w:\mathcal X\to Z:\;
\pi_w=\hat z_\mu^\#\circ \phi_w,\ \phi_w:\mathcal X\to \Xi_\mu^{\mathrm{synthetic}},\ w\in\Omega
\right\}.
\end{equation}

A natural question to ask is under which conditions optimal decisions can be represented by scenarios at all. As we shall see in the next section, our standing assumptions are sufficient for guaranteeing this.

\section{Optimal policy realizability via cost vector and demand vector learning}
\label{sec:optimal-policy-realizability}
\label{sec:realizability}

\begin{theorem}[Optimal policy realizability via cost vector and demand vector learning]
\label{thm:realizability}
Fix a context $x\in\mathcal X$ and a scenario $\xi\in\Xi$. Let $\mu \in [0,\infty)$ denote any admissible choice of log-barrier smoothing parameter, including $0$.

\textbf{Case 1: Learning the cost $q$}.
There exists a cost vector
$q_\mu^\#(x,\xi)\in\R^{n_2}$ which realizes an optimal first-stage decision, in the sense that
\begin{equation}
\label{eq:q-decision-optimal}
\argmin_{z\in\R^{n_1}}
\hat G_\mu\!\left(
z,
\left(W(\xi),h(\xi),q_\mu^\#(x,\xi)\right)
\right)
\cap
\argmin_{z\in\R^{n_1}}
\E_{\boldsymbol{\xi}\sim P(\cdot\mid x)}
\!\left[\hat G_\mu(z,\boldsymbol{\xi})\right]
\neq\emptyset.
\end{equation}
\textbf{Case 2: Learning the demand $h$}.
If the recourse is deterministic $\forall \xi\in\Xi$, $W(\xi)=W_{\det}$, there likewise exists a decision optimal demand vector
$h_{\mu,\delta}^\#(x)\in\R^{m_2}$ such that
\begin{equation}
\label{eq:h-decision-optimal}
\argmin_{z\in\R^{n_1}}
\hat G_\mu\!\left(
z,
\left(W_{\det},h_{\mu,\delta}^\#(x),q_\delta^{\max}\right)
\right)
\cap
\argmin_{z\in\R^{n_1}}
\E_{\boldsymbol{\xi}\sim P(\cdot\mid x)}
\!\left[\hat G_\mu(z,\boldsymbol{\xi})\right]
\neq\emptyset.
\end{equation}
where we define
\begin{equation}
\label{eq:qmaxdelta}
q_\delta^{\max}[i]:=\max_{\xi'\in\Xi}\left(q(\xi')[i]+\delta\right),
\qquad i=1,\ldots,n_2,\quad \delta>0.
\end{equation}
\end{theorem}

\begin{remark}[Synthetic scenarios]
\label{rem:synthetic-scenarios}
While the theorem may seem magic at first, we emphasize that it does not guarantee that the decision optimal scenario will exist in the support of the distribution of $\xi$. As mentioned, the scenario is ``synthetic'', in the sense that it need not correspond to a physically realistic real-life outcome. It is primarily a decision parameter.
\end{remark}

\begin{proposition}[Constructive expressions]
\label{prop:constructive}
The optimal scenario components $q_\mu^\#(x,\xi)$ and $h_{\mu,\delta}^\#(x)$ have closed form expressions in terms of $z^\#(x)$, which can be evaluated at computational cost at most equivalent to solving a single-scenario problem. See Appendix~\ref{app:realizability-proof}.
\end{proposition}

The above theorem shows that a scenario generator can realize an optimal decision policy, even under the restriction of generating only one scenario, or even a single component of a single scenario. Decision focused scenario generation is hence equally expressive as direct decision-policy learning.

By Proposition~\ref{prop:constructive}, constructive formulas are available for decision-optimal scenarios in terms of optimal decisions. Of course, one can likewise construct an optimal decision from an optimal scenario by solving a single-scenario problem. Hence, if we know the value of the decision-optimal $q$, $h$, or $z$, we can infer the value of the others. Learning any of them is hence equivalent in terms of interpretability and realizability.

Additionally, Proposition~\ref{prop:constructive} and Theorem~\ref{thm:realizability} show that the components you learn do not have to even correspond to the structure of the randomness in the stochastic program. For example, nothing prevents using a method for learning $q$, even if $q$ is deterministic in the original problem.

Generally, the choice of learning $q$ or $h$ via decision focused scenario generation, or $z$ directly through decision policy learning, should be made entirely based on the ease of the learning itself, as determined by dimensionality and complexity of the constraint sets of $\Xi$ and $Z$ respectively.

\section{Decision focused learning via log-barrier smoothing}
\label{sec:dfl-logbarrier}
\label{sec:barrier}
\label{sec:backprop-logbarriers}

For learning decision-optimized scenario policies, deep learning is a natural approach.
To train a neural net with back-propagation, we would need to differentiate the scenario-parametrized decision regret
\begin{equation}
\label{eq:scenario-parametrized-erm}
\min_w\ \frac1N\sum_{i=1}^N
\hat G_\mu\!\left(\hat z_\mu^\#(\phi_w(x_i)),\xi_i\right),
\qquad
\phi_w:\mathcal X\to \Xi_\mu^{\mathrm{synthetic}}.
\end{equation}
Formally, for a fixed index $i$, write $\hat\xi_i(w):=\phi_w(x_i)$ and $z_i(\hat\xi_i):=\hat z_\mu^\#(\hat\xi_i)$.
Differentiating with respect to the generated scenario gives
\begin{equation}
\label{eq:scenario-chain-rule}
\frac{\partial}{\partial \hat\xi_i}\hat G_\mu(z_i(\hat\xi_i),\xi_i)
=
\frac{\partial \hat G_\mu(z_i(\hat\xi_i),\xi_i)}{\partial z}
\frac{\partial \hat z_\mu^\#(\hat\xi_i)}{\partial \hat\xi_i}.
\end{equation}
The first term $\partial \hat G_\mu/\partial z$ can be handled via subgradients due to the convexity of $\hat G_\mu$.
The second factor presents more of an issue. In particular, the solution is generally piecewise constant in the cost $q$. Log-barrier smoothing resolves this pathology.

\begin{proposition}[Closed form derivatives for log-barriers]
\label{prop:logbarrier-derivatives}
Fix $\mu>0$, $\tilde\mu\geq 0$, and a realized scenario $\tilde\xi\in\Xi$. The scenario parametrized decision loss
\[
\mathcal L_{\mu,\tilde\mu}^{\tilde\xi}(\xi)
:=
\hat G_{\tilde\mu}\!\left(\hat z_\mu^\#(\xi),\tilde\xi\right),
\qquad
\xi=(W,h,q),
\]
is smooth in $q$ and $h$, and $\partial \hat z_\mu^\#/\partial q$ and $\partial \hat z_\mu^\#/\partial h$ have closed form derivatives in terms of $\hat z_\mu^\#(W,h,q)$ and its dual. See Appendix~\ref{app:closed-form-derivatives}.
\end{proposition}

In addition to enabling back-propagation, log-barrier smoothing yields a guarantee that the scenario-parametrized decision loss has no local minima that are not global:

\begin{theorem}[No vanishing gradients]
\label{thm:no-vanishing-gradients}
Fix $x \in \mathcal X$, $\mu > 0$. Fix an arbitrary recourse matrix from the support, $W_f:=W(\tilde\xi)$ for some $\tilde\xi\in\Xi$.
The cost-and-demand parametrized decision loss
\[
L_\mu(q,h,W_f,x)
:=
\E_{\boldsymbol{\xi}\sim P(\cdot\mid x)}
\!\left[
\hat G_\mu\!\left(\hat z_\mu^\#(q,h,W_f),\boldsymbol{\xi}\right)
\right]
\]
is componentwise invex in $q$ and $h$.
For a pair $(q^\#,h^\#)\in\R^{n_2}\times\R^{m_2}$, the following conditions are equivalent:
\begin{equation}
\label{eq:no-vg-equivalences}
\resizebox{\textwidth}{!}{$
\nabla_q L_\mu(q^\#,h^\#,W_f,x)=0
\;\Longleftrightarrow\;
\nabla_h L_\mu(q^\#,h^\#,W_f,x)=0
\;\Longleftrightarrow\;
\hat z_\mu^\#(q^\#,h^\#,W_f)=\argmin_{z\in\R^{n_1}}G_\mu(z,x).
$}
\end{equation}
\end{theorem}

The theorem implies that the loss landscape over the neural net cost and demand output is ``trap free''.
Training may still get stuck in local minima, as is generally the risk in deep learning, but only due to non-convexities in the net itself, not in the loss function.
The result is unexpected and has a non-trivial proof. In particular, it shows that log-barrier smoothing is theoretically distinct from quadratic smoothing.

In connection with Theorem~\ref{thm:realizability}, the equivalence in \eqref{eq:no-vg-equivalences} also provides an interesting certificate for the existence of an optimal $h$ in cases with random recourse $W$. A demand vector that solves the demand-parametrized decision problem exists if and only if that demand vector realizes an optimal scenario.

\section{An interior point inspired method for decision focused scenario generation}
\label{sec:interior-point-method}
\label{sec:method}
\label{sec:nnipm}

The previous section shows that log-barrier smoothing can enable decision focused scenario learning via backpropagation.
But we face a dilemma when setting the smoothing parameter $\mu$.
When $\mu$ is large, gradients will be well-behaved and stable, but the solution of the problem will be distorted.
Classical interior point methods proceed by first solving the problem for a sufficiently large $\mu$, and then progressively decreasing it while tracing the optimal solution via Newton's method.
Below, we present a novel method, applying the same philosophy for decision focused learning.

In an initial stage, we set $\mu$ to a relatively high value, to make training easy and stable.
In a subsequent homotopy stage, we progressively decrease $\mu$ while continuing training, maintaining low training loss as the problem approaches the non-smooth version.
When $\mu$ reaches a sufficiently small target value $\mu_f:=\mu_S$, we begin a final fine-tuning stage where we train directly on the unsmoothed decision loss, solving the partially unsmoothed empirical risk minimization problem
\begin{equation}
\label{eq:nnipm-finetune}
\min_w\
\frac1N\sum_{i=1}^N
\hat G\!\left(\hat z_{\mu_f}^\#(\phi_w(x_i)),\xi_i\right).
\end{equation}
Back-propagation through the unsmoothed loss $\hat G$ is handled via subgradients; see Appendix~\ref{app:closed-form-derivatives}.

\begin{algorithm}[H]
\caption{Neural Net Interior Point Method for scenario generation.}
\label{alg:nnipm-training}
\begin{algorithmic}[1]
\Statex \textbf{Input:} training data $\mathcal D=\{(x_i,\xi_i)\}_{i=1}^N$; scenario generator $\phi_w$; decreasing barrier schedule $\mu_1>\cdots>\mu_S>0$.
\Statex \textbf{Output:} trained parameters $w$.
\State Initialize $w$.
\For{$s=1$ \textbf{to} $S$}
    \State Set $\mu\gets\mu_s$.
    \For{training iterations}
        \State Sample minibatch $\mathcal B\subseteq\mathcal D$.
        \State Generate synthetic scenarios $\hat\xi_i\gets\phi_w(x_i)$ for $(x_i,\xi_i)\in\mathcal B$.
        \State Solve $z_i\gets\hat z_\mu^\#(\hat\xi_i)$ for each generated scenario.
        \State Update $w$ by backpropagating the minibatch loss $\frac1{|\mathcal B|}\sum_{(x_i,\xi_i)\in\mathcal B}\hat G_\mu(z_i,\xi_i)$.
    \EndFor
\EndFor
\For{fine-tuning iterations}
    \State Sample minibatch $\mathcal B\subseteq\mathcal D$.
    \State Generate synthetic scenarios $\hat\xi_i\gets\phi_w(x_i)$ for $(x_i,\xi_i)\in\mathcal B$.
    \State Solve $z_i\gets\hat z_{\mu_S}^\#(\hat\xi_i)$ for each generated scenario.
    \State Update $w$ by backpropagating the true downstream loss $\frac1{|\mathcal B|}\sum_{(x_i,\xi_i)\in\mathcal B}\hat G(z_i,\xi_i)$.
\EndFor
\State \textbf{return} $w$.
\end{algorithmic}
\end{algorithm}

We emphasize that this method can be applied for decision focused learning of the cost $q$, the demand $h$, or both, regardless of which components of $W$, $T$, $h$, and $q$ are random in the underlying stochastic program. Assumptions made previously in the paper are required for guarantees, but not for applicability. Empirical validation will determine which learning strategy works best for problems in practice.


\section{Experiments}
\label{sec:experiments}

The experiments are intended as a modest validation of the theory and algorithm. We use three contextual two-stage stochastic LP benchmarks chosen to isolate the three second-stage uncertainty regimes relevant to our theory. The resource-allocation benchmark is inherited from the application-driven pointwise forecasting literature and represents the fixed-recourse, fixed-cost case where uncertainty enters through the second-stage right-hand side. The transshipment benchmark is based on the multilocation transshipment literature and on random-cost stochastic-programming instances, and tests the mixed case where both demand and second-stage costs vary. Finally, the random-yield benchmark is a controlled synthetic stress test for random recourse matrices, motivated by production settings with uncertain yields. Together, these instances are not intended as an exhaustive empirical benchmark suite; rather, they validate the central modeling claim that the learned synthetic component can be chosen for computational and statistical convenience, rather than forced to match the physical source of uncertainty.

We focus on two questions. First, can the proposed log-barrier single-scenario generator be trained stably and achieve reasonable decision quality? Second, how much does the choice of learned scenario component matter in finite-sample training?

For each context $x$, the learned generator outputs a single synthetic scenario. The associated single-scenario problem is solved to obtain a first-stage decision $z(x)$, which is then evaluated on held-out scenarios. We report relative regret against a per-context sample-average oracle:
\[
\mathrm{RelRegret}
=
100\cdot
\frac{
\widehat G(z(x),x)-\widehat G(z^\star(x),x)
}{
|\widehat G(z^\star(x),x)|
}.
\]
Lower is better. The numbers below should be read as exploratory evidence that the theoretical construction leads to a trainable method, not as a comprehensive benchmark study.

\begin{table}[t]
\centering
\small
\caption{Effect of the learned scenario component. Entries are relative regret in percent. The true randomness indicates which part of the real stochastic problem varies; the learned component indicates which vector the neural generator outputs.}
\label{tab:component-choice}
\begin{tabular}{lccc}
\hline
Benchmark & True randomness & Learn $h$ & Learn $q$ \\
\hline
Resource allocation & $h$ & \textbf{8.07} & 24.31 \\
Transshipment & $h,q$ & \textbf{2.06} & 26.95 \\
Random yield & $W$ & 8.40 & \textbf{7.08} \\
\hline
\end{tabular}
\end{table}

The method trains stably on all instances in Table~\ref{tab:component-choice} and obtains single-digit relative regret when the learned component is chosen appropriately. The results also show that the learned component can have a large and non-obvious effect. In the transshipment problem with both random demand and random cost, learning demand gives $2.06\%$ regret, whereas learning the cost vector gives $26.95\%$ regret under the same evaluation protocol. Thus, although Theorem~\ref{thm:realizability} gives an existence result for decision-optimal single-scenario representations, finite-sample learning still strongly depends on which scenario component is parameterized.

The random-yield experiment illustrates the complementary phenomenon. In this instance the real uncertainty is in the recourse matrix $W$, which would be expensive and structurally delicate to learn directly. Instead, following the logic of Theorem~\ref{thm:realizability} and Proposition~\ref{prop:constructive}, we keep $W$ fixed in the generated scenario and learn a positive second-stage cost vector $q$. This achieves $7.08\%$ relative regret, slightly better than learning $h$ and substantially below the naive alternatives observed in our internal runs. This supports the main modeling point: the synthetic scenario need not be a literal forecast of the physical randomness in order to induce good first-stage decisions.

Overall, the experiments support the intended qualitative conclusions. Single-scenario log-barrier training works in principle on contextual two-stage stochastic LPs. More importantly, the learned scenario component should be treated as a design choice: in some problems learning demand is far better, while in others learning a cost vector can avoid the need to learn recourse matrices altogether.

\section{Conclusion}
\label{sec:conclusion}

This paper develops a theory of decision-focused single-scenario learning for contextual two-stage stochastic programming. The main message is that a learned scenario need not be an accurate sample from the conditional distribution in order to be decision optimal. Under the stated well-posedness assumptions, optimal stochastic-programming decisions can be induced by carefully chosen synthetic single-scenario problems.

This perspective separates decision parametrization from distribution forecasting. It explains why cost, demand, and direct decision learning can be equivalent from the standpoint of the first-stage decision they induce, while still differing substantially in trainability, dimensionality, and feasibility. The constructive formulas in Proposition~\ref{prop:constructive} make this equivalence explicit: decision-optimal synthetic costs and demands can be written in terms of optimal decisions and averaged dual information. The counterexamples in the appendix also show that the assumptions are not merely technical; fixing the wrong structural component can destroy representability.

The second part of the paper studies trainability. Log-barrier smoothing turns the scenario-to-decision map into a differentiable object with closed-form implicit derivatives. More importantly, the induced decision-focused loss has a componentwise invex structure in the learned cost and demand variables.

The experiments are intended as a validation of this theory rather than as a broad empirical benchmark.

Several directions remain open. Natural extensions include risk-averse objectives, richer scenario restrictions that preserve interpretability, and broader convex recourse models such as conic or semidefinite programs. A systematic scaling study on larger industrial instances is also needed to understand when the single-scenario theory translates into computational advantage in practice.

\bibliographystyle{plainnat}
\bibliography{references}

\appendix

\begin{center}
{\Large\bfseries Core appendix}
\end{center}

\section{Preliminary Definitions and Results}
\label{app:preliminary-definitions-results}

Throughout, we will repeatedly make use of the dual of the second stage problem, as defined by a scenario
\(\xi \in \Xi^{\mathrm{synthetic}}\) and a first stage decision \(z \in Z\).
By classical duality theory, it is equivalent in value to the second-stage primal problem.

\[
Q_\mu(z,\xi)
=
\mu n_2(1-\log\mu)
+
\max_{\lambda\in\mathbb R^{m_2}}
\left\{
\lambda^\top(h(\xi)-Tz)
+
\mu\sum_{k=1}^{n_2}
\log\bigl(q(\xi)-W(\xi)^\top\lambda\bigr)_k
:
q(\xi)-W(\xi)^\top\lambda>0
\right\}.
\]

As in the primal case, for \(\mu > 0\), the solution is guaranteed to be unique.

\[
\hat\lambda_\mu(z,\xi)
:=
\operatorname*{arg\,max}_{\lambda\in\mathbb R^{m_2}}
\left\{
\lambda^\top(h(\xi)-Tz)
+
\mu\sum_{k=1}^{n_2}
\log\bigl(q(\xi)-W(\xi)^\top\lambda\bigr)_k
:
q(\xi)-W(\xi)^\top\lambda>0
\right\}.
\]

\paragraph{KKT conditions.}

Classical theory provides a useful optimality condition in terms of the dual multipliers. Here \(N_Z(z)\) denotes the normal cone to \(Z\) at \(z\).

Likewise, for a single scenario \(\xi\),
\(z\in\arg\min \widehat G_\mu(\cdot,\xi)\) if and only if

\[
z\in Z\cap \mathbb R^{n_1}_{++},
\qquad
0
\in
c-\mu z^{-1}
-
T^\top\hat\lambda_\mu(z,\xi)
+
N_Z(z).
\]

For any \(x\) and any \(\mu>0\),
\(z\in \arg\min G_\mu(\cdot,x)\) if and only if

\[
z\in Z\cap \mathbb R^{n_1}_{++},
\qquad
0
\in
c-\mu z^{-1}
-
T^\top
\mathbb E_{\xi\sim P(\cdot\mid x)}
\left[
\hat\lambda_\mu(z,\xi)
\right]
+
N_Z(z).
\]

In particular, for \(\mu>0\),

\[
\nabla_z Q_\mu(z,\xi)
=
-T^\top\hat\lambda_\mu(z,\xi).
\]

\section{Optimal Decision Representability}
\label{app:realizability-proof}


\subsection{Closed-Form Expressions for Optimal Decision Inducing Demand and Cost}

This subsection records the closed-form scenario components used in Proposition~\ref{prop:constructive} and Theorem~\ref{thm:realizability}. The dual variables are those defined in Appendix~\ref{app:preliminary-definitions-results}.

\paragraph{Decision-optimal cost vector.}
For a reference scenario $\tilde\xi=(\widetilde W,\widetilde h,\widetilde q)$, let $\widetilde\lambda_\mu^x$ be any second-stage dual optimizer for $\tilde\xi$ at $z_\mu^\#(x)$. Define
\begin{equation}
\label{eq:core-q-construct}
q_\mu^{\#,\tilde\xi}(x)
:=
\widetilde q+
\widetilde W^\top
\left(
\lambda_\mu^{\#,x}-\widetilde\lambda_\mu^x
\right).
\end{equation}
This is the cost-vector construction used in Theorem~\ref{thm:realizability}.

\paragraph{Decision-optimal demand vector.}
Assume that the recourse matrix is deterministic, $W(\xi)=W$ for all scenarios, and fix $\delta>0$. Define $q_\delta^{\max}$ as in \eqref{eq:qmaxdelta}. For $\mu>0$, set
\begin{equation}
\label{eq:core-h-residual}
r_\mu^{\#,x,\delta}
:=
q_\delta^{\max}-W^\top\lambda_\mu^{\#,x},
\end{equation}
\begin{equation}
\label{eq:core-y-construct}
y_\mu^{\#,\delta}(x)
:=
\mu\left(r_\mu^{\#,x,\delta}\right)^{-1},
\end{equation}
and
\begin{equation}
\label{eq:core-h-construct}
h_\mu^{\#,\delta}(x)
:=
Tz_\mu^\#(x)+Wy_\mu^{\#,\delta}(x).
\end{equation}
For $\mu=0$, set
\begin{equation}
\label{eq:core-h-construct-zero}
h_0^{\#,\delta}(x):=Tz_0^\#(x).
\end{equation}
These are the right-hand-side constructions used in Theorem~\ref{thm:realizability}.

\subsection{Proof of Theorem 1}

\begin{proof}[Proof of Theorem~\ref{thm:realizability}]
Fix a context $x$ and a smoothing parameter $\mu\geq 0$. For $\mu>0$, define the averaged second-stage dual multiplier by
\[
\lambda_\mu^{\#,x}
:=
\E_{\boldsymbol\xi\sim P(\cdot\mid x)}
\left[
\hat\lambda_\mu\!\left(z_\mu^\#(x),\boldsymbol\xi\right)
\right].
\]
For $\mu=0$, let $\lambda_0^{\#,x}$ be any averaged second-stage dual multiplier at $z_0^\#(x)$ satisfying the stochastic KKT condition.

Define
\[
\beta_\mu(z)
:=
\begin{cases}
\mu z^{-1}, & \mu>0,\\
0, & \mu=0.
\end{cases}
\]
Here $N_Z(z)$ denotes the normal cone to $Z$ at $z$. A point $z$ is optimal for a single-scenario problem with scenario $\xi=(W,h,q)$ if there exists a second-stage dual optimizer $\lambda$ for that scenario at $z$ such that
\[
0\in c-\beta_\mu(z)-T^\top\lambda+N_Z(z).
\]
Similarly, since $z_\mu^\#(x)$ is optimal for the stochastic problem, there exists an averaged multiplier $\lambda_\mu^{\#,x}$ such that
\begin{equation}
\label{eq:realizability-proof-stochastic-kkt}
0\in c-\beta_\mu(z_\mu^\#(x))-T^\top\lambda_\mu^{\#,x}+N_Z(z_\mu^\#(x)).
\end{equation}

We first prove the $q$-construction. Fix a reference scenario
\[
\tilde\xi=(\widetilde W,\widetilde h,\widetilde q).
\]
Choose any second-stage dual optimizer $\widetilde\lambda_\mu$ of the reference scenario at $z_\mu^\#(x)$. Define
\[
q_\mu^{\#,\tilde\xi}(x)
:=
\widetilde q+\widetilde W^\top(\lambda_\mu^{\#,x}-\widetilde\lambda_\mu).
\]
We claim that $\lambda_\mu^{\#,x}$ is a second-stage dual optimizer for the synthetic scenario
\[
(\widetilde W,\widetilde h,q_\mu^{\#,\tilde\xi}(x))
\]
at the same first-stage point $z_\mu^\#(x)$.

Indeed, for any $\lambda$, define
\[
\bar\lambda
:=
\lambda-\lambda_\mu^{\#,x}+\widetilde\lambda_\mu.
\]
Then
\[
q_\mu^{\#,\tilde\xi}(x)-\widetilde W^\top\lambda
=
\widetilde q-\widetilde W^\top\bar\lambda.
\]
Thus $\lambda$ is feasible for the synthetic dual if and only if $\bar\lambda$ is feasible for the reference dual; this holds both for the strict feasible set when $\mu>0$ and the closed feasible set when $\mu=0$. Moreover, the synthetic dual objective at $\lambda$ equals the reference dual objective at $\bar\lambda$ plus the constant
\[
(\lambda_\mu^{\#,x}-\widetilde\lambda_\mu)^\top(\widetilde h-Tz_\mu^\#(x)).
\]
Therefore the maximizers are translated by $\lambda_\mu^{\#,x}-\widetilde\lambda_\mu$. Since $\widetilde\lambda_\mu$ is optimal for the reference scenario, $\lambda_\mu^{\#,x}$ is optimal for the synthetic scenario.

Combining this second-stage optimality with the stochastic first-stage KKT condition \eqref{eq:realizability-proof-stochastic-kkt}, we get
\[
0\in c-\beta_\mu(z_\mu^\#(x))-T^\top\lambda_\mu^{\#,x}+N_Z(z_\mu^\#(x)),
\]
where $\lambda_\mu^{\#,x}$ is now a valid dual optimizer for the synthetic single-scenario problem. Hence $z_\mu^\#(x)$ satisfies the single-scenario KKT condition for
\[
(\widetilde W,\widetilde h,q_\mu^{\#,\tilde\xi}(x)).
\]
By convexity, $z_\mu^\#(x)$ is optimal for that single-scenario problem. This proves $q$-realizability for both $\mu>0$ and $\mu=0$.

Now consider the $h$-construction. Assume $W$ is deterministic across scenarios. Fix $\delta>0$ and define $q_\delta^{\max}$ componentwise by
\[
q_\delta^{\max}[j]
=
\max_{\xi\in\Xi}q(\xi)[j]+\delta.
\]
For $\mu>0$, define
\[
r_\mu^{\#,x,\delta}
:=
q_\delta^{\max}-W^\top\lambda_\mu^{\#,x}.
\]
This vector is strictly positive componentwise: writing $\lambda_{\mu,\xi}$ for a supported dual optimizer at $z_\mu^\#(x)$, $q_\delta^{\max}$ dominates every supported cost vector by at least $\delta$, and every supported dual multiplier satisfies $q(\xi)-W^\top\lambda_{\mu,\xi}>0$. Define
\[
y_{\mu,\delta}^\#(x)
:=
\mu\left(r_\mu^{\#,x,\delta}\right)^{-1}
\]
and
\[
h_{\mu,\delta}^\#(x)
:=
Tz_\mu^\#(x)+Wy_{\mu,\delta}^\#(x).
\]
Then, by construction,
\[
Wy_{\mu,\delta}^\#(x)
=
h_{\mu,\delta}^\#(x)-Tz_\mu^\#(x),
\]
and
\[
q_\delta^{\max}-W^\top\lambda_\mu^{\#,x}
=
\mu\left(y_{\mu,\delta}^\#(x)\right)^{-1}.
\]
Thus $\lambda_\mu^{\#,x}$ is the second-stage dual optimizer for the synthetic scenario
\[
(W,h_{\mu,\delta}^\#(x),q_\delta^{\max})
\]
at $z_\mu^\#(x)$. Combining this with the stochastic KKT condition \eqref{eq:realizability-proof-stochastic-kkt}, $z_\mu^\#(x)$ satisfies the single-scenario KKT condition. Hence, by convexity, $z_\mu^\#(x)$ is optimal for the synthetic single-scenario problem. This proves $h$-realizability for $\mu>0$.

Finally, consider $\mu=0$ for the $h$-construction. Set
\[
h_{0,\delta}^\#(x)
:=
Tz_0^\#(x).
\]
The dual feasible set
\[
\{\lambda:\;W^\top\lambda\leq q_\delta^{\max}\}
\]
contains every supported second-stage dual feasible set, because $q_\delta^{\max}$ dominates every $q(\xi)$. Therefore it contains the averaged stochastic multiplier $\lambda_0^{\#,x}$. Since
\[
h_{0,\delta}^\#(x)-Tz_0^\#(x)=0,
\]
the unsmoothed second-stage dual objective is identically zero over the feasible dual polytope. Hence every feasible dual point is optimal, in particular $\lambda_0^{\#,x}$. Combining this with the stochastic KKT condition
\[
0\in c-T^\top\lambda_0^{\#,x}+N_Z(z_0^\#(x)),
\]
we conclude that $z_0^\#(x)$ satisfies the single-scenario KKT condition for
\[
(W,h_{0,\delta}^\#(x),q_\delta^{\max}).
\]
Therefore $z_0^\#(x)$ is optimal for that synthetic single-scenario problem. This proves the $h$-realizability claim for $\mu=0$.
\end{proof}


\section{Log-Barrier Smoothing}
\label{app:log-barrier-smoothing}


\subsection{Closed-Form Derivatives and Proof of Theorem 2}
\label{app:closed-form-derivatives}
\label{app:no-vanishing-gradients-proof}

Fix $\mu>0$ and $(W,h,q)\in\Xi_\mu^{\mathrm{synthetic}}$. For the rest of this subsection, write $\lambda(\xi,z):=\hat\lambda_\mu(z,\xi)$ for the second-stage dual optimizer defined in Appendix~\ref{app:preliminary-definitions-results}. Set
\[
        \hat z:=\hat z_\mu^\#(W,h,q),
        \qquad
        \hat\lambda:=\lambda((W,h,q),\hat z).
\]
All derivatives are taken on the relative interior of the domain. Vector powers are componentwise.

For
\[
        L_\mu(q,h,W,x):=G_\mu(\hat z_\mu^\#(W,h,q),x),
\]
the chain rule gives
\[
        \nabla_h L_\mu(q,h,W,x)
        =
        \left(\frac{\partial \hat z}{\partial h}\right)^\top
        \nabla_zG_\mu(\hat z,x),
        \qquad
        \nabla_q L_\mu(q,h,W,x)
        =
        \left(\frac{\partial \hat z}{\partial q}\right)^\top
        \nabla_zG_\mu(\hat z,x).
\]
By the envelope theorem,
\[
        \nabla_z \widehat G_\mu(z,\xi)
        =
        c-\mu z^{-1}-T^\top\lambda(\xi,z)
        \quad \bmod\operatorname{range}(A^\top),
\]
and hence
\[
        \nabla_z G_\mu(z,x)
        =
        c-\mu z^{-1}
        -
        T^\top \E_{\xi\sim P(\cdot\mid x)}[\lambda(\xi,z)]
        \quad \bmod\operatorname{range}(A^\top).
\]

\begin{lemma}
Let $B\in\mathbb R^{n\times m}$ have full column rank and let
\[
        \lambda(a,r)
        :=
        \argmax_{\lambda:\ a-B\lambda>0}
        \left\{
        r^\top\lambda
        +
        \mu\sum_{i=1}^n \log(a_i-(B\lambda)_i)
        \right\}.
\]
At a solution, set
\[
        D:=\operatorname{Diag}\bigl((a-B\lambda(a,r))^{-2}\bigr).
\]
Then
\[
        d\lambda
        =
        (\mu B^\top D B)^{-1}
        \bigl(dr+\mu B^\top D\,da\bigr).
\]
Equivalently,
\[
        \frac{\partial\lambda}{\partial r}
        =
        (\mu B^\top D B)^{-1},
        \qquad
        \frac{\partial\lambda}{\partial a}
        =
        (\mu B^\top D B)^{-1}\mu B^\top D .
\]
\end{lemma}

\begin{proof}
The first-order equation is
\[
        r-\mu B^\top(a-B\lambda)^{-1}=0.
\]
Differentiating gives
\[
        dr+\mu B^\top D\,da-\mu B^\top D B\,d\lambda=0.
\]
Since $B$ has full column rank and $D\succ0$, $\mu B^\top D B$ is nonsingular.
\end{proof}

Applying the lemma with $B=W^\top$, $a=q$, and $r=h-Tz$ gives
\[
        d\lambda
        =
        (\mu WDW^\top)^{-1}
        \bigl(dh-T\,dz+\mu WD\,dq\bigr),
        \qquad
        D:=\operatorname{Diag}\bigl((q-W^\top\hat\lambda)^{-2}\bigr).
\]
In particular, at fixed $(W,h,q)$,
\[
        \frac{\partial\lambda}{\partial z}
        =
        -(\mu WDW^\top)^{-1}T.
\]

The KKT equations for $\hat z$ are
\[
        c-\mu \hat z^{-1}-T^\top\hat\lambda+A^\top\nu=0,
        \qquad
        A\hat z=b.
\]
Let
\[
        D_z:=\operatorname{Diag}(\hat z^2),
        \qquad
        P
        :=
        \frac1\mu
        \left(
        D_z
        -
        D_zA^\top(AD_zA^\top)^{-1}AD_z
        \right).
\]
Differentiating the first-stage KKT equations gives
\[
        d\hat z=P T^\top d\hat\lambda.
\]
Combining this identity with the preceding expression for $d\lambda$ yields
\[
        S\,d\hat\lambda=dh+\mu WD\,dq,
        \qquad
        S:=TPT^\top+\mu WDW^\top.
\]
Therefore
\[
        \frac{\partial\hat\lambda}{\partial h}=S^{-1},
        \qquad
        \frac{\partial\hat\lambda}{\partial q}=S^{-1}\mu WD,
\]
and
\[
        \boxed{
        \frac{\partial \hat z_\mu^\#}{\partial h}
        =
        PT^\top S^{-1}}
        ,
        \qquad
        \boxed{
        \frac{\partial \hat z_\mu^\#}{\partial q}
        =
        PT^\top S^{-1}\mu WD}.
\]
Thus, for any representative
\[
        g\in\nabla_zG_\mu(\hat z,x),
\]
the loss gradients are
\[
        \boxed{
        \nabla_h L_\mu(q,h,W,x)
        =
        S^{-1}TPg},
        \qquad
        \boxed{
        \nabla_q L_\mu(q,h,W,x)
        =
        \mu DW^\top S^{-1}TPg}.
\]

\paragraph{Proof of Theorem 2.}
Fix $x\in \mathcal X$, $\mu>0$, and $W_f$ from the support. Write $W:=W_f$ and
\[
        \hat z:=\hat z_\mu^\#(W,h,q),
        \qquad
        g\in\nabla_zG_\mu(\hat z,x).
\]

First assume
\[
        d\in\ker A,
        \quad
        Td=0
        \quad\Longrightarrow\quad
        d=0.
\]
Since $P=P^\top$, $\operatorname{range}(P)=\ker A$, $P\succ0$ on $\ker A$, and $W$ has full row rank, the closed-form gradients give
\[
        \nabla_h L_\mu(q,h,W,x)=0
        \iff
        TPg=0,
\]
and
\[
        \nabla_q L_\mu(q,h,W,x)=0
        \iff
        \mu DW^\top S^{-1}TPg=0
        \iff
        TPg=0.
\]
The rank condition gives
\[
        TPg=0
        \iff
        Pg=0.
\]
Finally,
\[
        Pg=0
        \iff
        d^\top g=0
        \quad\text{for every }d\in\ker A.
\]
Hence, under the rank condition,
\[
        \nabla_h L_\mu(q,h,W,x)=0
        \iff
        \nabla_q L_\mu(q,h,W,x)=0
        \iff
        d^\top\nabla_zG_\mu(\hat z,x)=0
        \quad\forall d\in\ker A.
\]

It remains to remove the rank condition. Set
\[
        u:=Tz,
        \qquad
        Z_T:=T(Z),
\]
and define
\[
        \phi_\mu(u)
        :=
        \argmin_{\substack{z\in Z\cap\mathbb R^{n_1}_{++}\\ Tz=u}}
        \left\{
        c^\top z-\mu\sum_{j=1}^{n_1}\log z_j
        \right\}.
\]
Also define
\[
        F_\mu(u)
        :=
        c^\top\phi_\mu(u)
        -
        \mu\sum_{j=1}^{n_1}\log(\phi_\mu(u)_j),
\]
and
\[
        Q_\mu^T(u,\xi)
        :=
        \min_{\substack{y>0\\ W(\xi)y=h(\xi)-u}}
        \left\{
        q(\xi)^\top y-\mu\sum_{k=1}^{n_2}\log y_k
        \right\}.
\]
Then
\[
        \bar G_\mu(u,x)
        :=
        F_\mu(u)
        +
        \E_{\xi\sim P(\cdot\mid x)}
        [Q_\mu^T(u,\xi)]
\]
satisfies
\[
        G_\mu(\phi_\mu(u),x)=\bar G_\mu(u,x).
\]
Likewise, for the synthetic scenario $(W,h,q)$,
\[
        \hat u_\mu(h,q)
        :=
        \argmin_{u\in Z_T}
        \left\{
        F_\mu(u)+Q_\mu^T(u,(W,h,q))
        \right\},
\]
and
\[
        \hat z_\mu^\#(W,h,q)=\phi_\mu(\hat u_\mu(h,q)),
        \qquad
        L_\mu(q,h,W,x)=\bar G_\mu(\hat u_\mu(h,q),x).
\]

In the reduced problem the technology map is the identity on
\[
        \operatorname{lin}(Z_T-Z_T)=T(\ker A),
\]
so the rank-condition case applies. Hence
\[
\begin{aligned}
        \nabla_h L_\mu(q,h,W,x)=0
        &\iff
        \nabla_q L_\mu(q,h,W,x)=0
        \\
        &\iff
        v^\top\nabla_u\bar G_\mu(\hat u_\mu(h,q),x)=0
        \quad\forall v\in\operatorname{lin}(Z_T-Z_T).
\end{aligned}
\]
By the envelope theorem for the fiber minimization,
\[
        d^\top\nabla_zG_\mu(\phi_\mu(u),x)
        =
        (Td)^\top\nabla_u\bar G_\mu(u,x),
        \qquad d\in\ker A.
\]
Therefore, in the original variables,
\[
        \nabla_h L_\mu(q,h,W,x)=0
        \iff
        \nabla_q L_\mu(q,h,W,x)=0
        \iff
        d^\top\nabla_zG_\mu(\hat z,x)=0
        \quad\forall d\in\ker A.
\]
Since $G_\mu(\cdot,x)$ is convex on $Z$ and has the unique minimizer $z_\mu^\#(x)$, the last condition is equivalent to
\[
        \hat z_\mu^\#(W,h,q)=z_\mu^\#(x).
\]
The rest of the theorem statement follows immediately.

\section{Tightness of Theorem 1}

Theorem~\ref{thm:realizability} makes several assumptions. As we shall see in this section, the conclusions of the theorem cannot generally be made to hold while relaxing these assumptions.

\citet{wallace2000decision} presents elementary examples of problems where no scenario in the support of the noise yields an optimal first-stage decision. \citet{homemdemello2024forecasting} go further, showing that restrictions on the admissible scenario components can prevent decision-optimal single-scenario representability. Hence, fully synthetic scenarios are generally required, though there may be problems where restrictions can be imposed without losing optimal decision representability.

Below, we present additional counterexamples showing that our assumptions of deterministic technology matrix $T$, and deterministic recourse for learning $h$, cannot be relaxed in general.

\paragraph{Counterexample showing the necessity of deterministic $T$.}
\label{app:deterministic-T-counterexample}

The assumption that $T$ is deterministic cannot be dropped if the synthetic scenario is required to keep a fixed support value of $T$.

Consider first the inequality description
\[
0\leq u_1\leq 1,
\qquad
0\leq u_2\leq 1.
\]
In the standard equality form used in the paper, write
\[
z=(u_1,u_2,s_1,s_2)\geq 0,
\qquad
Az=b,
\]
where
\[
A=
\begin{pmatrix}
1&0&1&0\\
0&1&0&1
\end{pmatrix},
\qquad
b=(1,1).
\]
Thus $u_1+s_1=1$ and $u_2+s_2=1$. Let
\[
c=(1,1,0,0).
\]
There are two equiprobable scenarios. The second-stage problem is
\[
Q_i(z)=
\min q_i^\top y
\quad
\text{subject to } W_i y=h_i-T_i z,\quad y\geq 0.
\]
Let
\[
W_1=W_2=
\begin{pmatrix}
1&-1
\end{pmatrix},
\qquad
h_1=h_2=1,
\]
\[
T_1=
\begin{pmatrix}
1&0&0&0
\end{pmatrix},
\qquad
T_2=
\begin{pmatrix}
0&1&0&0
\end{pmatrix},
\]
and
\[
q_1=(1,1),
\qquad
q_2=(4,4).
\]
The log-barrier version of this standard-form problem is well defined: the first-stage feasible set has strictly positive points, every recourse equation admits strictly positive $y$ for every feasible $z$, and the only nonzero recourse recession directions are of the form $(t,t)$, on which the recourse costs are strictly positive.

For scenario 1,
\[
y^+-y^-=1-u_1.
\]
Since $0\leq u_1\leq 1$, this gives
\[
Q_1(z)=1-u_1.
\]
Similarly,
\[
Q_2(z)=4(1-u_2).
\]
Therefore the stochastic LP objective is
\[
\begin{aligned}
G(z)
&=c^\top z+\frac{1}{2}Q_1(z)+\frac{1}{2}Q_2(z)\\
&=u_1+u_2+\frac{1}{2}(1-u_1)+\frac{1}{2}4(1-u_2)\\
&=\frac{5}{2}+\frac{1}{2}u_1-u_2.
\end{aligned}
\]
Hence the unique stochastic optimizer is
\[
u^\#=(0,1),
\]
or, in standard-form variables,
\[
z^\#=(0,1,1,0).
\]
Now fix the first support technology
\[
T=T_1=
\begin{pmatrix}
1&0&0&0
\end{pmatrix}.
\]
For any choice of $W$, $h$, and $q$, the synthetic recourse value can depend on $z$ only through
\[
Tz=u_1.
\]
Thus the synthetic single-scenario objective has the form
\[
u_1+u_2+\psi(u_1),
\]
with $\psi$ determined by the chosen $W$, $h$, and $q$. In particular, the recourse term is independent of $u_2$.

If a feasible synthetic solution has $u_2>0$, replacing $u_2$ by $0$ and $s_2$ by $1$ preserves first-stage feasibility and leaves the recourse value unchanged, while strictly decreasing the first-stage cost. Therefore every optimizer of any well-defined synthetic single-scenario LP with fixed $T=T_1$ must satisfy
\[
u_2=0.
\]
This is incompatible with the stochastic optimizer, which satisfies $u_2^\#=1$. Hence no choice of $W$, $h$, and $q$ can reproduce the stochastic optimizer when $T$ is fixed to $T_1$.

This shows that deterministic $T$, or an equivalent row-space condition, is genuinely needed.

\paragraph{Counterexample showing the need for deterministic $W$ when learning $h$.}
\label{app:deterministic-W-counterexample}

The $h$-only construction also fails if $W$ is not deterministic.

Use the same first-stage set. In inequality form,
\[
0\leq u_1\leq 1,
\qquad
0\leq u_2\leq 1.
\]
In standard equality form,
\[
z=(u_1,u_2,s_1,s_2)\geq 0,
\qquad
Az=b,
\]
with
\[
A=
\begin{pmatrix}
1&0&1&0\\
0&1&0&1
\end{pmatrix},
\qquad
b=(1,1).
\]
Let
\[
c=(-3,-3,0,0),
\]
and let
\[
u^\#=\left(\frac{1}{2},\frac{1}{2}\right),
\]
so that
\[
z^\#=\left(\frac{1}{2},\frac{1}{2},\frac{1}{2},\frac{1}{2}\right).
\]
Let the technology matrix be deterministic:
\[
T=
\begin{pmatrix}
-1&0&0&0\\
0&-1&0&0
\end{pmatrix}.
\]
Thus
\[
h-Tz=h+(u_1,u_2).
\]
There are two equiprobable scenarios. Let
\[
h_1=h_2=\left(-\frac{1}{2},-\frac{1}{2}\right),
\]
and
\[
q_1=q_2=q=(1,6,1,1).
\]
Write
\[
y=(a_1,a_2,b_1,b_2)\geq 0.
\]
In scenario 1, set
\[
W_1=
\begin{pmatrix}
1&0&-1&0\\
0&1&0&-1
\end{pmatrix}.
\]
In scenario 2, let
\[
P=
\begin{pmatrix}
0&1\\
1&0
\end{pmatrix},
\]
and set
\[
W_2=\begin{pmatrix}P&-P\end{pmatrix},
\]
that is,
\[
W_2=
\begin{pmatrix}
0&1&0&-1\\
1&0&-1&0
\end{pmatrix}.
\]
The log-barrier version of this standard-form problem is well defined: the first-stage feasible set has strictly positive points, every recourse equation admits strictly positive $y$ for every feasible $z$, and all nonzero recourse recession directions have strictly positive cost under $q$.

For scenario 1,
\[
a_1-b_1=u_1-\frac{1}{2},
\qquad
a_2-b_2=u_2-\frac{1}{2}.
\]
Therefore
\[
\begin{aligned}
Q_1(z)
&=\max\left\{u_1-\frac{1}{2},0\right\}
+6\max\left\{u_2-\frac{1}{2},0\right\}\\
&\quad
+\max\left\{\frac{1}{2}-u_1,0\right\}
+\max\left\{\frac{1}{2}-u_2,0\right\}.
\end{aligned}
\]
For scenario 2,
\[
P(a-b)=u-u^\#,
\]
or equivalently
\[
a-b=P(u-u^\#).
\]
Hence
\[
\begin{aligned}
Q_2(z)
&=\max\left\{u_2-\frac{1}{2},0\right\}
+6\max\left\{u_1-\frac{1}{2},0\right\}\\
&\quad
+\max\left\{\frac{1}{2}-u_2,0\right\}
+\max\left\{\frac{1}{2}-u_1,0\right\}.
\end{aligned}
\]
Averaging the two scenarios gives, up to an additive constant,
\[
\begin{aligned}
G(z)
&=c^\top z+\frac{1}{2}Q_1(z)+\frac{1}{2}Q_2(z)\\
&=\mathrm{constant}
+\sum_{j=1}^{2}\left[
\frac{1}{2}\max\left\{u_j-\frac{1}{2},0\right\}
+4\max\left\{\frac{1}{2}-u_j,0\right\}
\right].
\end{aligned}
\]
This function is minimized uniquely at
\[
u_1=u_2=\frac{1}{2}.
\]
Thus the unique stochastic optimizer is
\[
z^\#=\left(\frac{1}{2},\frac{1}{2},\frac{1}{2},\frac{1}{2}\right).
\]

Now try to learn only $h$ while fixing $W$ to one support value. Let
\[
q^{\max}=q+(1,1,1,1)=(2,7,2,2).
\]
First fix
\[
W=W_1.
\]
For any choice of $h$, the dual feasible set of the synthetic recourse problem is
\[
\begin{aligned}
D_1
&=\{\lambda\in\R^2:\; W_1^\top\lambda\leq q^{\max}\}\\
&=[-2,2]\times[-2,7].
\end{aligned}
\]
Since
\[
T=\begin{pmatrix}-I&0\end{pmatrix},
\]
a recourse subgradient with respect to the first two first-stage coordinates is $\lambda$. Because $z^\#$ is strictly positive, the nonnegativity constraints are inactive at $z^\#$. If $z^\#$ were optimal for the synthetic single-scenario problem, the first-stage KKT conditions would require some $\lambda\in D_1$ and some equality multiplier $\alpha=(\alpha_1,\alpha_2)$ satisfying
\[
0=c+(\lambda_1,\lambda_2,0,0)+A^\top\alpha.
\]
The slack coordinates give
\[
\alpha_1=\alpha_2=0.
\]
The first two coordinates then force
\[
\lambda_1=3,
\qquad
\lambda_2=3.
\]
Thus optimality of $z^\#$ would require
\[
(3,3)\in D_1.
\]
But
\[
D_1=[-2,2]\times[-2,7],
\]
so this is impossible.

Now fix
\[
W=W_2.
\]
Then
\[
\begin{aligned}
D_2
&=\{\lambda\in\R^2:\; W_2^\top\lambda\leq q^{\max}\}\\
&=[-2,7]\times[-2,2].
\end{aligned}
\]
The same KKT argument again forces $\lambda=(3,3)$, but
\[
(3,3)\notin D_2.
\]
Therefore no choice of $h$ can make $z^\#$ optimal when $W$ is fixed to $W_1$ or $W_2$.

This shows that learning $h$ can shift the second-stage right-hand side, but it cannot enlarge the dual feasible set $\{\lambda:\; W^\top\lambda\leq q\}$. When $W$ varies across scenarios, the stochastic optimum may require a dual slope that is unavailable under any fixed support value of $W$.

\section{Details on Experiments and Benchmarks Against Standard Methods}
\label{app:details-experiments-benchmarks}

This appendix records the experimental protocol, benchmark definitions, implementation details, hardware, and baseline comparisons used in the empirical study. All methods are evaluated as contextual two-stage stochastic linear-programming policies: a policy observes a context, returns a first-stage decision, and is scored by downstream stochastic-programming cost.

\subsection{Experimental Protocol}

Each benchmark uses the same data generating protocol. The training set contains $100$ contexts with one observed scenario per context. The test set contains $30$ independently generated contexts. For each test context, we generate $500$ conditional scenarios and evaluate policies over five batches of $100$ scenarios. Reported values are averages over the resulting $30\times 5=150$ test evaluations. The test contexts and scenarios are generated independently from the training data.

For each test context and evaluation batch, we compute both the candidate policy cost and an oracle SAA cost on the same batch. For a policy $\pi$, regret is reported as
\[
\widehat{\mathrm{Regret}}(\pi)
=
\frac{1}{30\cdot 5}
\sum_{i=1}^{30}
\sum_{b=1}^{5}
\left[
\widehat V_{i,b}(\pi(x_i))-\widehat V_{i,b}^{\star}
\right],
\]
where $\widehat V_{i,b}^{\star}$ is the oracle SAA value for test context $i$ and batch $b$. Relative regret divides each gap by the corresponding oracle value before averaging. Lower values are better.


\subsection{Benchmarks}

\begin{table}[h]
\centering
\small
\begin{tabular}{lcccc}
\toprule
Benchmark & Random component & $n_1$ & $n_2$ & $m_2$ \\
\midrule
Resource allocation & $h$ & $20$ & $680$ & $50$ \\
Transshipment & $h,q$ & $7$ & $77$ & $35$ \\
Random yield & $W$ & $5$ & $20$ & $5$ \\
\bottomrule
\end{tabular}
\caption{Benchmark families and equality-form dimensions.}
\label{tab:app-benchmark-dimensions}
\end{table}

\paragraph{Problem descriptions and data generation.}
We give a more explicit description of the three benchmark families emphasized in the main text. All three are contextual two-stage stochastic linear programs. A context is observed before the first-stage decision is chosen, and uncertainty is then realized through one or more second-stage data components. Two of the benchmarks are inherited from prior contextual or stochastic-programming benchmark papers; the third is a controlled synthetic stress test targeting the random-recourse-matrix regime that prior single-scenario methods largely do not cover.

\paragraph{Resource allocation.}
The resource-allocation instance is included as a direct point of contact with the recent application-driven pointwise forecasting literature. It is a contextual two-stage stochastic LP with uncertainty entering only through the second-stage right-hand side, matching the fixed-recourse and fixed-cost setting studied by \citet{homemdemello2024forecasting}. The instance also connects to data-driven SAA with covariate information \citep{kannan2025datadriven}. This makes it a natural baseline regime for testing whether the proposed log-barrier single-scenario training procedure can recover useful decision policies in a setting where single-scenario representability is already known to be favorable.

In the data generator, each context is three-dimensional. Contexts are generated as folded Gaussian vectors with a randomly generated correlation matrix. Conditional demand is then produced by combining a contextual signal with additive noise; in the reported experiments the noise level is $\sigma=5$ and the contextual exponent is $p=2$. The base instance has $20$ resources and $30$ demand classes, so the generated demand vector has length $30$. In equality form, this demand vector is embedded into the second-stage right-hand side by prepending $20$ zero entries for resource-balance rows, giving $m_2=50$ second-stage equality rows. The full equality-form instance has $n_1=20$ first-stage variables and $n_2=680$ second-stage variables. The random component is $h$, while the recourse matrix $W$, second-stage cost vector $q$, and technology matrix $T$ are fixed.

\paragraph{Transshipment.}
The transshipment benchmark is included to test the method in a random-cost stochastic-programming regime. The base problem is a multilocation transshipment model from the inventory and supply-chain literature, where post-demand lateral transshipments are represented through a network-flow recourse problem \citep{herer2006multilocation}. We use it in the spirit of the random-cost stochastic-programming experiments of \citet{gangammanavar2021stochastic}, who introduced random second-stage objective coefficients into classical two-stage stochastic LP test instances. This benchmark is therefore qualitatively different from the resource-allocation instance: uncertainty affects both the second-stage right-hand side and second-stage objective vector.

The data generator uses three-dimensional contexts centered around four fixed context regimes. Conditional scenarios are generated by perturbing both the second-stage right-hand side and the second-stage cost vector. The underlying transshipment instance is read from SMPS data and has $n_1=7$ first-stage variables, $n_2=77$ second-stage variables, and $m_2=35$ second-stage equality rows. In the mixed version used here, the contextual generator varies $7$ stochastic right-hand-side values and $7$ stochastic recourse-cost values. The main comparison stresses the choice of learned scenario component: the real stochastic problem contains both demand and cost variation, but we may still choose to learn only one of the two.

\paragraph{Random yield.}
The random-yield instance is a controlled synthetic stress test for random recourse matrices. Unlike the previous two benchmarks, it is not meant to reproduce a canonical named test instance. Instead, it isolates a source of uncertainty that is central to the paper's theoretical motivation: the physical effectiveness of second-stage actions may vary across scenarios, causing uncertainty to enter the recourse matrix $W$. Such random-recourse structures are studied in stochastic programming \citep{fan2024decision,jiang2019quantitative}, and arise naturally in production-planning settings with uncertain yields \citep{kazemi2011production}. The benchmark tests whether a synthetic vector-valued scenario, especially a learned cost vector $q$, can induce good first-stage decisions without requiring the neural generator to output a full recourse matrix.

This benchmark is included because directly learning a recourse matrix can be expensive and structurally delicate. The generator therefore creates scenarios in which the true conditional distribution changes $W$, while the learned synthetic scenario is restricted to vector-valued data such as $h$ or $q$. In the small instance used here, the equality-form dimensions are $n_1=5$, $n_2=20$, and $m_2=5$. The construction uses $r=5$ products, $a=10$ activities, and a support size $K_{\mathrm{support}}=5$, giving $n_2=a+2r=20$. The random recourse matrix has the block form $W=[Y\ I\ -I]$, where $Y$ is a scenario-dependent yield matrix. The base cost vector has length $20$: ten activity costs, five positive-slack costs, and five negative-slack costs. This benchmark directly tests the modeling message of Theorem~\ref{thm:realizability} and Proposition~\ref{prop:constructive}: a synthetic scenario need not reproduce the physical source of uncertainty in order to induce a good first-stage decision.

\subsection{Methods and Hyperparameters}

The baseline set is
\[
\text{SAA},\quad
\text{OLS-CE},\quad
\text{ER-SAA},\quad
\text{NN}.
\]
SAA is context-free and solves one sample-average approximation using all training scenarios. OLS-CE fits an ordinary least-squares map from context to the relevant scenario component and solves the induced one-scenario problem. ER-SAA fits the same linear map but evaluates an SAA formed by adding empirical residuals to the predicted component. NN uses the same neural architecture as IP-SG, but is trained by mean-squared prediction error rather than decision loss.

All MLP-based policies use three hidden ReLU layers of width $128$ followed by a benchmark-specific output layer. Neural scenario generators are trained with Adam, learning rate $10^{-3}$, minibatch size $1$, and one generated scenario per context. Training uses $130$ epochs: $20$ epochs for the first smoothing stage and $10$ epochs for each later stage. The log-barrier schedule is
\[
\mu\in
\{1.0,\ 0.8,\ 0.6,\ 0.4,\ 0.2,\ 0.1,\ 0.08,\ 0.06,\ 0.04,\ 0.02,\ 0.01\}.
\]
During annealing, $(\mu_{\mathrm{in}},\mu_{\mathrm{ref}})=(\mu_s,\mu_s)$. The final fine-tuning stage uses $(\mu_{\mathrm{in}},\mu_{\mathrm{ref}})=(0.01,0)$.

All experiments are implemented in Julia. Neural-network training uses Flux. Smoothed log-barrier problems are solved with Ipopt, while unsmoothed LPs and oracle SAA problems are solved with HiGHS. Neural-network computations use Float64 precision.

\begin{table}[h]
\centering
\small
\begin{tabular}{ll}
\toprule
Item & Value \\
\midrule
Machine & MacBook Pro, model identifier Mac15,10 \\
Chip & Apple M3 Max \\
CPU & 14 cores total, 10 performance and 4 efficiency \\
GPU & Integrated Apple M3 Max GPU, 30 cores, Metal 3 \\
Memory & 36 GB unified memory \\
Operating system & macOS 15.6, build 24G84 \\
Kernel & Darwin 24.6.0 \\
\bottomrule
\end{tabular}
\caption{Hardware and operating-system details.}
\label{tab:app-hardware}
\end{table}

\clearpage
\subsection{Benchmark Results}

\begin{table}[H]
\centering
\scriptsize
\setlength{\tabcolsep}{2.5pt}
\resizebox{\textwidth}{!}{%
\begin{tabular}{lrrrrrr}
\toprule
Benchmark & SAA & ER & OLS & NN & IP-SG-$q$ & IP-SG-$h$ \\
\midrule
\texttt{random\_yield} & 23.62\% & 46.81\% & 29.14\% & 26.07\% & 7.08\% & 8.40\% \\
\texttt{resource\_allocation} & 23.54\% & 6.27\% & 8.52\% & 11.73\% & 24.31\% & 8.07\% \\
\texttt{transshipment} & 13.22\% & 4.08\% & 6.02\% & 6.99\% & 26.95\% & 2.06\% \\
\bottomrule
\end{tabular}%
}
\caption{Benchmark comparison.}
\label{tab:app-dfl-overlap-comparison}
\end{table}

The comparison is intended to contextualize the IP-SG results rather than to establish a comprehensive leaderboard.
Against the standard baselines, IP-SG is strongest on transshipment, where IP-SG-$h$ gives the lowest regret in the table.
On resource allocation, IP-SG-$h$ is competitive with the predictive baselines, though not the best overall.
On random yield, the cost-vector IP-SG variant is the strongest vector-valued generator and substantially improves over the non-scenario-generation baselines, despite the physical uncertainty being in the recourse matrix.
Across the board, IP-SG outperforms naive baselines when the learned component is chosen appropriately, but can fail otherwise.
This pattern matches the main empirical message of the paper: the learned synthetic component is a modeling choice, and its best choice need not coincide with the physical source of randomness.

\clearpage
\input{checklist}

\end{document}
